\section{A direct polynomial in twenty-seven unrestricted integers}
\label{sec:integer27}

The preceding nine-variable construction minimizes the number of nonnegative
unknowns among the constructions supplied here, but it has a very large degree.
We now give a direct signed-coordinate alternative. Combining the weighted
three-root elimination of package~09 with the signed-core argument of
package~14 reduces that direct family's unrestricted-integer count from
$29$ to $27$. Two versions have exact degrees $1004$ and $1116$, with
$137$ and $136$ arithmetic operations, respectively. The reduction below is
an ordinary mathematical proof; the new specialized circuits have not been
formalized in Lean or Rocq.

Throughout this section, $(\beta,K,M_C,M_F)$ is the complete four-port T4
coefficient tuple defined earlier, with $\beta=B-1$ and with $M_F$ the
\emph{shifted} mask. The proof uses the established native compiler interface
and its canonical padding by independent powers of five. The complete older
package-06 tuple is an alternative binding with the same properties; individual
ports from different tuples must not be combined.

\subsection{The polynomial and its domain}

Temporarily retain seven positive coordinates
\[
 J,F,\alpha,z,s,\eta,\zeta>0
\]
and six unrestricted integer coordinates
\[
 h,r,T,u,\gamma,t\in\mathbb Z.
\]
Set $i=2r+1$ and $w=u^2$, and use the following polynomial abbreviations:
\begin{align*}
 q&=\beta J+1,& X&=wq,& Y&=sq^3,& E&=XY,& k&=\eta+\zeta,\\
 c&=kY+\eta,& a&=E+Y,& H&=4a+3,& \Delta&=a^2+H.
\end{align*}
The main root, packed index, and three outer factors are
\begin{align*}
 D_{\rm main}&=X+ac+\gamma H,\\
 R&=((q-1)(q-F)+\alpha)(q^2-1)+(M_C+qM_F)J,\\
 N_{\rm main}&=D_{\rm main}^2-\Delta c^2,\\
 N_k&=k-hE-R,\\
 N_t&=(K+w)(q-F-\alpha)+q-F-z(q-1).
\end{align*}
For the eliminated roots define
\begin{align*}
 L_0&=EkY,& A&=1+L_0(L_0+k),& I&=ic^2,\\
 D&=1+\Delta I^2,& Q&=\Delta(D-1),& d&=Q-1,\\
 v&=cT,& z_0&=c+RD,& W_0&=Dv^2+z_0^2,\\
 e&=QW_0-1,& B_0&=2Qvz_0.
\end{align*}
The two admissible weights are
\[
 M_{\rm full}=D+v^2,
 \qquad M_{\rm low}=kI^2+T^2.
\]
For either choice $M$ put
\begin{align}
 C&=d(t^2+M^2A)-e,\nonumber\\
 U&=C^2+DB_0^2-4d^2t^2M^2A,\nonumber\\
 \mathcal R&=U^2-D(2CB_0)^2,\label{eq:i27-resultant}\\
 P_{27}&=N_{\rm main}N_kN_t(1+\mathcal R)-1.
 \label{eq:i27-polynomial}
\end{align}
All displayed symbols are substituted polynomials, not further quantified
unknowns. The alternative four-equation system is
\begin{equation}
 N_{\rm main}=N_k=N_t=1,\qquad \mathcal R=0.
 \label{eq:i27-system}
\end{equation}

To obtain literally $27$ unrestricted unknowns, replace each of the seven
positive coordinates $x$ by an independent block
\begin{equation}
 x=\operatorname{Pos}(a_x,b_x,c_x)
   =a_x^2+a_x+b_x^2+c_x^2+1.
 \label{eq:i27-pos}
\end{equation}
This is positive for every integer triple and represents every positive
integer. Indeed, $4x-3$ is positive and congruent to $1$ modulo $4$, so
Legendre's three-square theorem represents it as a sum of three squares.
Exactly one of those squares is odd and two are even; writing the three
roots as $2a_x+1,2b_x,2c_x$ gives \eqref{eq:i27-pos}. Thus the total is
$7\cdot3+6=27$, with no extra variable for a sign or a denominator.

\paragraph{Claim.}
For either weight, \eqref{eq:i27-polynomial} after
\eqref{eq:i27-pos}, and likewise \eqref{eq:i27-system}, has an integer
solution if and only if K215 halts on the blank tape. Consequently the
absence of integer solutions is equivalent in ZFC to
$\operatorname{Con}(\mathsf{ZFC})$.

\subsection{Unit extraction and bounds before restoring roots}

At a zero of $P_{27}$, all four factors are integer units. Since
$\Delta\equiv0$ or $3\pmod4$, the main norm cannot be $-1$.
Moreover $\mathcal R\equiv U^2\pmod4$, so $1+\mathcal R$ cannot be
$-1$. Therefore
\begin{equation}
 N_{\rm main}=1,\quad \mathcal R=0,\quad
 N_k=\epsilon,\quad N_t=\nu,\quad
 \epsilon,\nu\in\{-1,1\},\quad\epsilon\nu=1.
 \label{eq:i27-units}
\end{equation}
The system is the special case $\epsilon=\nu=1$.

Initially $X\ge0$, while $Y,k,c,a,H,\Delta$ are positive. If $u=0$,
then $X=E=0$, $a=Y$, and
\[
 N_{\rm main}
   =H(2ac\gamma+\gamma^2H-c^2).
\]
This cannot equal $1$ because $H=4Y+3>1$. Thus $u\ne0$, $w\ge1$,
and $X\ge q\ge16$. The inherited weak-transport argument and exact
mask bounds give
\begin{equation}
 q-F-\alpha\ge0,\qquad
 q(q^2-1)<R<q^4-2,\qquad E\ge q^4>R+2,
 \qquad a>R+2.
 \label{eq:i27-outer}
\end{equation}
These bounds use only the transport/index units and the coefficient
interface, not any of the deleted roots. A negative $h$ would give
$k=R+\epsilon+hE<0$, so $h\ge0$. Also $Y\ge q^3\ge4096$ and
$c<k(Y+1)$.

\subsection{The main norm already forces a positive index quotient}

This additional observation is needed to use the lower weight without
assuming the deleted first norm. Let $A_0=a+2$, so
$\Delta=A_0^2-1$. Write the standard Pell sequences as
\[
 \chi_{A_0}(n)+\psi_{A_0}(n)\sqrt{A_0^2-1}
   =(A_0+\sqrt{A_0^2-1})^n.
\]
The main norm gives
$c=\psi_{A_0}(p)$ and $|D_{\rm main}|=\chi_{A_0}(p)$ for $p\ge1$.
Since $c>1$, $p\ne1$.

If $p=2$ and $D_{\rm main}>0$, then
$D_{\rm main}-ac=H+4$. Its defining congruence would give
$X\equiv4\pmod H$, impossible because $16\le X<H$. If $p=2$ and
$D_{\rm main}<0$, direct expansion gives
$4(D_{\rm main}-ac)\equiv-1\pmod H$. This would give
$4X\equiv-1\pmod H$, impossible because $0<4X+1<H$.
Thus $p\ge3$ and
\[
 c\ge\psi_{A_0}(3)=4A_0^2-1.
\]
If $k\le E$, then
\[
 c<k(Y+1)\le E(Y+1)\le E(E+1)<4(E+2)^2-1<4A_0^2-1,
\]
a contradiction. Hence $k>E$. In view of \eqref{eq:i27-outer},
$h=0$ would imply $k=R+\epsilon<E$. We have therefore proved
\begin{equation}
 h\ge1,\qquad k>E,\qquad k>R
 \label{eq:i27-hpositive}
\end{equation}
before restoring the first Pell root.

\subsection{Strict domination for both weights}

We establish
\begin{equation}
 AM^2>4W_0,\qquad dM^2A>2QW_0.
 \label{eq:i27-domination}
\end{equation}
The second follows from the first because $Q\ge2$ and
$d=Q-1\ge Q/2$.

For the full weight, \eqref{eq:i27-hpositive} implies
\[
 c^2+R^2<k^2((Y+1)^2+1)\le5k^2Y^2,
 \qquad A>L_0^2=k^2X^2Y^4>16(c^2+R^2).
\]
Here $X\ge16$ and $Y\ge1$ give the last strict inequality. Since
$M^2\ge4Dv^2$, $M^2\ge D^2$ and
$z_0^2\le2D^2(c^2+R^2)$,
\[
 4W_0\le M^2\bigl(1+8(c^2+R^2)\bigr)<AM^2.
\]

For the lower weight, $A_0<2E$, $c<2kY$ and $L_0=kEY$ give
\begin{equation}
 A>\frac{A_0^2c^2}{16},\qquad
 A>\frac{A_0^4Y^2}{16}>32A_0^4.
 \label{eq:i27-low-estimates}
\end{equation}
The second uses $k>E$, already proved, and $Y\ge4096$.
Furthermore $I^2=i^2c^4>c$, so
\[
 D<A_0^2I^2,
 \qquad z_0=c+RD<2kA_0^2I^2.
\]
The inequalities $M\ge kI^2$ and $M^2\ge4kI^2T^2$, together with
$k>R>3q\ge48$ and the first estimate in
\eqref{eq:i27-low-estimates}, imply $AM^2>8Dv^2$; when $T=0$ this is
immediate. The second estimate gives $AM^2>8z_0^2$. Adding and dividing
by two proves \eqref{eq:i27-domination}. In particular, for every
integer $t$,
\begin{equation}
 C=d t^2+dM^2A-(QW_0-1)>d t^2+QW_0+1>0.
 \label{eq:i27-Cpositive}
\end{equation}

\subsection{Excluding the two new exceptional branches}

The original positive-variable elimination assumed $t,T>0$. Here they
are unrestricted, so their possible zeros require proof.

\paragraph{The branch $t=0$.}
At $t=0$, the resultant is
\[
 \mathcal R=(C^2-DB_0^2)^2.
\]
But $DB_0^2=4Q^2Dv^2z_0^2\le Q^2W_0^2$, whereas
\eqref{eq:i27-Cpositive} gives $C>QW_0+1$. Thus $\mathcal R>0$,
contradicting \eqref{eq:i27-units}. Hence $t\ne0$.

\paragraph{The branch $T=0$.}
First observe that
\begin{equation}
 2<z_0<Q-1.
 \label{eq:i27-zgap}
\end{equation}
The lower bound is immediate, and, because $\Delta>R$,
\begin{align*}
 Q-z_0
   &=(\Delta-R)D-\Delta-c\\
   &\ge D-\Delta-c
     =1+\Delta(i^2c^4-1)-c>1.
\end{align*}
If $T=0$, then $v=B_0=0$, and $\mathcal R=0$ gives
$C^2=4d^2t^2M^2A$. Since $d,t,M$ are nonzero,
$\tau=C/(2dtM)$ is rational with square $A\in\mathbb Z$ and hence
is an integer. Put $y=t-M\tau$. The definition of $C$ gives
$dy^2=e=Qz_0^2-1$, so
\[
 (Q-1)(y^2-z_0^2)=z_0^2-1>0.
\]
Integrality implies $|y|\ge z_0+1$. Using
\eqref{eq:i27-zgap} then yields the contradiction
\[
 z_0^2-1\ge(Q-1)(2z_0+1)>z_0(2z_0+1)>z_0^2-1.
\]
Therefore $T\ne0$, and hence $vB_0\ne0$. Every denominator in the
following reconstruction has now been proved nonzero.

\subsection{Restoring the signed source equations}

Define the proof-only rational numbers
\[
 f=-\frac{U}{2CB_0},\qquad
 \tau=\frac{C+B_0f}{2dtM},\qquad
 y=t-M\tau.
\]
Equation $\mathcal R=0$ gives $f^2=D$. The definition of $U$ then
gives $(C+B_0f)^2=4d^2t^2M^2A$, hence $\tau^2=A$. A rational number
whose square is an integer is integral, as follows by reducing its
fraction to lowest terms. Thus $f,\tau,y$ are integers. Expansion gives
\[
 dy^2=e-B_0f=Q(vf-z_0)^2-1.
\]
Since $f^2=D$, the expression $vf-z_0$ is exactly
\[
 V=c(Tf-1)-Rf^2.
\]
We have restored all three normalized equations
\begin{equation}
 \tau^2-L_0(L_0+k)=1,
 \qquad f^2-\Delta i^2c^4=1,
 \qquad QV^2-(Q-1)y^2=1.
 \label{eq:i27-restored}
\end{equation}
Together with \eqref{eq:i27-units}, these are exactly the normalized
source product equations. The recovered $f,\tau,y$ may have signs.

We now reuse package~14's \emph{signed-core soundness argument}: its
seven positive ports are retained, the multiplier $i=2r+1$ is nonzero,
the square scale $w=u^2$ has already been proved positive, and all the
source equations have been restored. That argument classifies absolute
Pell roots, uses the two auxiliary congruences to recover the exact main
index $p=R$, excludes the negative main-root branch by its short residue
interval, and recovers the native masks and the blank tableau. It allows
$f,h,T,\tau,y,\gamma$ to be signed; it does not infer soundness by
evaluating the positive-domain theorem at a signed tuple. In particular,
it forces the two remaining unit signs to be positive. The result is a
blank halting run of K215. This imported argument is invoked only after
\eqref{eq:i27-restored}, never to establish the reconstruction denominators.

\subsection{Completeness, including the T4 coefficient binding}

The canonical compiler witness has $R\equiv3\pmod4$ and
$c=\psi_{a+2}(R)$ odd. As in package~14, choose the strong index
$m=cR$ and put
\[
 f=\chi_{a+2}(m),\qquad i=\frac{\psi_{a+2}(m)}{c^2}.
\]
The expansion of
$(\chi_{a+2}(R)+c\sqrt\Delta)^c$ shows
$c^2\mid\psi_{a+2}(m)$. Since $m$ is odd, both the numerator and
$c^2$ are odd, so $i$ is a positive odd integer. The standard odd quotient
identities at $R\equiv3\pmod4$ give the required positive auxiliary
witnesses and the exact polynomial $V$.

The canonical spatial and temporal dimensions may independently be
sufficiently large powers of five. Their product $N$ is odd. For T4,
$B=2^{d_0}$ with $d_0=5^{18}$ odd, so
\[
 q=2^{d_0N},\qquad w=2^{R-d_0N}
\]
is a positive square: $R-d_0N$ is positive and even. This is the precise
reason that package~14's square-scale completeness survives the T4
replacement. The whole tuple has the same compiler interface; no port is
changed in isolation. The older package-06 value $d_0=5^{67}$ has the same
parity property.

Choose integers $r,u$ with $i=2r+1$ and $w=u^2$. For either weight set
$t=M\tau+y$. The forward identities give $\mathcal R=0$, while the
three outer factors equal $1$. Finally use the surjectivity of
\eqref{eq:i27-pos} to choose the seven independent integer triples.
This proves completeness for the literal $27$-integer polynomial and
system. The stated equivalence with consistency now follows from the
earlier K215 calibration; it is not a proof of consistency itself.

\subsection{Exact degrees, circuit lengths, and evidence}

\begin{center}
\begin{tabular}{lrrrr}
\hline
Weight and form & Variables & Equations & Exact degree & Operations\\
\hline
Full product & $27$ in $\mathbb Z$ & $1$ & $1116$ & $136$\\
Lower product & $27$ in $\mathbb Z$ & $1$ & $1004$ & $137$\\
Full system & $27$ in $\mathbb Z$ & $4$ & $1064$ maximum & $130$\\
Lower system & $27$ in $\mathbb Z$ & $4$ & $952$ maximum & $131$\\
\hline
\end{tabular}
\end{center}

The source products have $84$ and $85$ operations. Each of seven
positivity substitutions costs seven, the odd multiplier costs two, and
the square scale costs one; the increment is therefore $52$. Every binary
addition, subtraction and multiplication is counted, including a product
by a fixed numeral. The lower product has $73$ multiplications and
$64$ additions or subtractions. Removing the six finalizer gates gives
the corresponding systems.

For an exact degree proof, assign weight two to the seven positive source
ports and $w$, and weight one to $i,h,T,\gamma,t$. The relevant degrees are
\[
\begin{array}{c|rrrrrrrrr}
\text{expression}&q&X&Y&E&k&c&a&\Delta&R\\ \hline
\text{degree}&2&4&8&12&2&10&12&24&8
\end{array}
\]
and
\[
\begin{array}{c|rrrrrrr}
\text{expression}&A&I&D&Q&v&z_0&e\\ \hline
\text{degree}&44&21&66&90&11&74&238.
\end{array}
\]
The outer factor degrees are $35,13,4$. For the lower weight, $M$ has
degree $44$, and the unique leading term of $C$ is $-e_{\rm top}$.
The unique leaders of $U$ and $\mathcal R$ are then $C_{\rm top}^2$
and $U_{\rm top}^2$, giving degrees $238,476,952$; the competing last
norm term has degree $892$. Hence the product degree is
$952+35+13+4=1004$. For the full weight, $M$ has degree $66$, and
$dM^2A$ dominates $C$. The successive degrees are $266,532,1064$, giving
the product degree $1116$.

There is no hidden cancellation of the outer leaders. They are respectively
\[
 8\gamma a_{\rm top}^{,2}c_{\rm top},\qquad
 -hE_{\rm top},\qquad
 w(q_{\rm top}-F_{\rm top}-\alpha_{\rm top})-z_{\rm top}q_{\rm top}.
\]
Each is nonzero, and their product with the resultant leader is nonzero.
The top parts of \eqref{eq:i27-pos}, $i=2r+1$, and $w=u^2$ preserve
nonvanishing: set the second and third variable of every positivity block
to zero, leaving an injective substitution into distinct square variables
and the nonzero scalar substitution $i\mapsto2r$. These degrees are exact
for every valid fixed $\beta>0$.

The supplied generator emits all four literal circuits, checks operation
counts and complete liveness, and compares both weights with an independent
formula on $96$ signed assignments each. It also checks $144$ isolated
signed root reconstructions, $6786$ finite instances of the gap inequality,
and two modular degree diagnostics per circuit. Both ordinary and optimized
Python runs passed. These are finite algebraic checks; the universal proof
is the argument above, and the inherited compiler and K215 calibration
remain explicit mathematical dependencies.
